% =============================================================================
% Discussion snippets for QSAR/QM noise robustness paper
% Paste into appropriate sections of the Discussion
% =============================================================================


% ---- (a) NGBoost distributional robustness ----
% Suggested location: Discussion, after presenting NGBoost results
% ----

The strong noise robustness of NGBoost can be understood through its
distributional learning framework~\citep{duan2020ngboost}. Unlike conventional
regression models that minimize mean squared error against (potentially noisy)
targets, NGBoost parameterizes the full conditional distribution $p(y \mid x)$
as a Normal distribution with learned mean $\mu(x)$ and variance $\sigma^2(x)$
via natural gradient boosting on the negative log-likelihood. This
parameterization provides an implicit mechanism for noise absorption: when a
training sample carries anomalously high label noise, the model can ``explain
away'' the discrepancy by increasing the local variance estimate $\sigma^2(x)$
rather than distorting the mean prediction $\mu(x)$ to fit the noisy target. In
effect, the variance head acts as a noise sink---observations that are difficult
to predict precisely are assigned high predictive variance, which down-weights
their influence on the mean prediction through the likelihood. This stands in
contrast to MSE-trained models, where every residual contributes equally to the
loss regardless of whether it reflects genuine structure or injected noise. Our
results are consistent with this mechanism: NGBoost configurations maintain
shallower degradation slopes across noise levels, suggesting that the
distributional output absorbs a substantial fraction of the injected noise into
the variance parameter rather than allowing it to corrupt point predictions.


% ---- (b) Tree-based outlier isolation ----
% Suggested location: Discussion, explaining RF/XGBoost robustness patterns
% ----

The inherent robustness of tree-based ensembles to label noise arises from
multiple complementary mechanisms rooted in the partition-and-aggregate
structure of random forests~\citep{breiman2001random}. First, recursive binary
splitting naturally isolates outliers---including samples with corrupted
labels---into their own terminal leaves containing few observations. Since
predictions are made by averaging within leaves, an outlier confined to a
small leaf has minimal influence on the majority of predictions across the
feature space. Second, bootstrap aggregation ensures that any individual noisy
sample appears in approximately 63.2\% of the constituent trees and, even in
those trees, only affects predictions for the local region defined by its leaf
partition. The ensemble averaging across hundreds of trees further dilutes the
influence of any single corrupted observation~\citep{mentch2016quantifying}.
These structural properties explain why the outlier noise strategy---which by
design corrupts only the $\sim$5\% of samples exceeding $|z| > 2.0$---has
negligible impact on tree-based model performance: the affected samples
represent a small fraction of the training set, are likely to be isolated in
sparse leaves during splitting, and their influence is averaged away across the
ensemble. This natural outlier containment, requiring no explicit robust loss
function or preprocessing, makes tree-based ensembles particularly well-suited
for molecular property prediction tasks where label noise may be concentrated in
a small subset of compounds.


% ---- (c) Mapping to Frenay & Verleysen noise taxonomy ----
% Suggested location: Discussion or Methods, connecting strategies to literature
% ----

Our six noise injection strategies map onto the taxonomy of label noise proposed
by~\citet{frenay2014classification}, providing a systematic framework for
interpreting robustness patterns. The \textit{legacy} strategy, which applies
additive Gaussian noise with constant variance $\sigma^2$ to all targets
regardless of feature values or target magnitude, corresponds to Noise
Completely At Random (NCAR): the corruption probability and magnitude are
independent of both $\mathbf{x}$ and $y$. The \textit{outlier},
\textit{quantile}, and \textit{threshold} strategies correspond to Noise Not At
Random (NNAR), where the corruption depends on the true target value $y$
(e.g., only samples with $|z| > 2.0$ or $y$ in the top/bottom quantiles are
corrupted), but not on the input features $\mathbf{x}$. Finally, the
\textit{heteroscedastic} and \textit{value-proportional} strategies implement
Noise At Random (NAR), where the noise magnitude scales with properties of the
target value itself---proportionally for value-proportional noise and via a
variance function $\sigma(y) = \sigma_{\text{base}} \cdot (1 + \alpha|y|)$ for
heteroscedastic noise. While Fr{\'e}nay and Verleysen's original taxonomy was
developed for classification, its extension to continuous-valued regression
targets is natural: the key distinction between NCAR, NNAR, and NAR lies in
the conditional dependence structure of the noise process, which applies
equally to categorical and continuous labels. Notably, systematic empirical
comparison of model robustness across all three noise categories remains
largely unexplored for regression tasks, representing a gap that our
experimental design directly addresses.


% ---- (d) Aleatoric vs epistemic uncertainty decomposition ----
% Suggested location: Discussion, interpreting uncertainty calibration results
% ----

The decomposition of predictive uncertainty into aleatoric and epistemic
components~\citep{kendall2017uncertainties} provides a principled lens for
understanding why certain models maintain calibrated uncertainty estimates under
noise injection while others do not. Aleatoric uncertainty captures irreducible
noise inherent in the data---measurement error, stochastic variation, or, in our
experimental setup, the artificially injected label noise. Epistemic uncertainty,
by contrast, reflects model ignorance due to insufficient data or
expressiveness, and can in principle be reduced by collecting more training
samples. A noise-robust model should attribute the injected noise primarily to
increased aleatoric uncertainty: correctly identifying the corruption as
irreducible data-level variation rather than a signal to be learned. In
contrast, a model that responds to noise injection by inflating epistemic
uncertainty is, in effect, becoming ``confused''---interpreting the noisy
labels as evidence that its own structure or parameterization is inadequate,
rather than recognizing the noise as an external perturbation. Our uncertainty
analysis tests this hypothesis directly: for distributional models such as
NGBoost and conformal prediction methods, we examine whether increasing
$\sigma$ leads to proportional growth in the aleatoric component while leaving
epistemic estimates relatively stable. Configurations that achieve this clean
separation are those best positioned for deployment in real-world molecular
property prediction, where label noise from experimental assays is endemic and
reliable uncertainty quantification is essential for downstream decision-making
in drug discovery pipelines.
